%===== FILE: HE2002_revision_notes.tex =====
\documentclass[12pt]{article}

%==================== PACKAGES ====================
\usepackage[margin=1in]{geometry}
\usepackage{amsmath,amssymb,amsthm}
\usepackage{enumitem}
\usepackage{bm}
\usepackage{hyperref}
\usepackage{fancyhdr}
\usepackage{xcolor}
\usepackage{tabularx}
\usepackage{array}
\usepackage{booktabs}
\usepackage{tikz}
\usepackage{pgfplots}
\pgfplotsset{compat=1.17}
\usepackage[most]{tcolorbox}
\tcbuselibrary{theorems,skins,breakable}
\usepackage{graphicx}
\usepackage{multicol}
\usepackage{draftwatermark}

%==================================================
% COURSE METADATA (EDIT THESE FOR EACH MODULE)
%==================================================
\newcommand{\CourseCode}{HE2002}
\newcommand{\CourseName}{Macroeconomics II}
\newcommand{\DocType}{Revision Notes}
\newcommand{\AcademicYear}{Academic Year 2025--2026}

%==================================================
% TCOLORBOX THEOREM STYLES
%==================================================

% ---------- DEFINITION: dark green, title in darker band ----------
\newtcbtheorem[number within=section]{definition}{Definition}%
{enhanced,
breakable,
colback=green!5!white,        % light interior
colframe=green!40!black,      % border
colbacktitle=green!50!black,  % dark title band
coltitle=white,               % title text colour
fonttitle=\bfseries,
title filled,                 % puts title inside the band
boxed title style={
size=small,
top=2pt, bottom=2pt,
left=6pt, right=6pt,
arc=1pt, outer arc=1pt
},
boxrule=0.9pt,
arc=1pt,
left=8pt, right=8pt, top=10pt, bottom=8pt,
before skip=12pt, after skip=12pt
}{def}

% ---------- THEOREM: blue, title in darker band ----------
\newtcbtheorem[number within=section]{theorem}{Theorem}%
{enhanced,
breakable,
colback=blue!3!white,
colframe=blue!40!black,
colbacktitle=blue!55!black,
coltitle=white,
fonttitle=\bfseries,
title filled,
boxed title style={
size=small,
top=2pt, bottom=2pt,
left=6pt, right=6pt,
arc=1pt, outer arc=1pt
},
boxrule=0.9pt,
arc=1pt,
left=8pt, right=8pt, top=10pt, bottom=8pt,
before skip=12pt, after skip=12pt
}{thm}

% ---------- ENVIRONMENTS USING TCOLORBOXENVIRONMENT ----------
\tcolorboxenvironment{identity}{
enhanced,
breakable,
colback=green!5!white,
colframe=green!60!black,
fonttitle=\bfseries,
boxrule=1pt,
arc=3pt,
left=8pt, right=8pt, top=8pt, bottom=8pt,
before skip=12pt, after skip=12pt
}

\tcolorboxenvironment{principle}{
enhanced,
breakable,
colback=red!5!white,
colframe=red!60!black,
fonttitle=\bfseries,
boxrule=1pt,
arc=3pt,
left=8pt, right=8pt, top=8pt, bottom=8pt,
before skip=12pt, after skip=12pt
}

%==================================================
% CUSTOM TCOLORBOX ENVIRONMENTS (WITH NAMES)
%==================================================
% Optional argument [1][] lets you override the default title if needed:
%   \begin{keyformula}[title={Key Formula: Demand Curve}] ... \end{keyformula}

\newtcolorbox{keyformula}[1][]{
enhanced,
breakable,
colback=blue!5!white,
colframe=blue!75!black,
title=Key Formula,
fonttitle=\bfseries,
boxrule=0.9pt,
arc=2pt,
left=8pt, right=8pt, top=8pt, bottom=8pt,
#1
}

% *** CHANGED TO BROWN AS REQUESTED ***
\newtcolorbox{examplebox}[1][]{
enhanced,
breakable,
colback=brown!5!white,
colframe=brown!75!black,
title=Worked Example,
fonttitle=\bfseries,
boxrule=0.9pt,
arc=2pt,
left=8pt, right=8pt, top=8pt, bottom=8pt,
#1
}

\newtcolorbox{commonmistake}[1][]{
enhanced,
breakable,
colback=red!5!white,
colframe=red!75!black,
title=Common Mistake,
fonttitle=\bfseries,
boxrule=0.9pt,
arc=2pt,
left=8pt, right=8pt, top=8pt, bottom=8pt,
#1
}

\newtcolorbox{policynote}[1][]{
enhanced,
breakable,
colback=yellow!10!white,
colframe=orange!75!black,
title=Policy Application,
fonttitle=\bfseries,
boxrule=0.9pt,
arc=2pt,
left=8pt, right=8pt, top=8pt, bottom=8pt,
#1
}

\newtcolorbox{examtip}[1][]{
enhanced,
breakable,
colback=purple!5!white,
colframe=purple!75!black,
title=Exam Focus,
fonttitle=\bfseries,
boxrule=0.9pt,
arc=2pt,
left=8pt, right=8pt, top=8pt, bottom=8pt,
#1
}

%==================================================
% TITLE / HEADER / FOOTER / WATERMARK
%==================================================

\title{%
\vspace{-0.5cm}
{\Large\bfseries \CourseCode\ \CourseName{} -- \DocType}
}

\author{Quantitative Research Society @ NTU}
\date{\AcademicYear}

\pagestyle{fancy}
\fancyhf{}
\fancyhead[L]{\CourseCode\ \CourseName{} -- \DocType}
\fancyhead[R]{\href{https://qrsntu.org}{QRS@NTU}}
\fancyfoot[C]{\thepage}
\renewcommand{\headrulewidth}{0.4pt}
\renewcommand{\footrulewidth}{0pt}

% Watermark
\SetWatermarkText{QRS@NTU — qrsntu.org}
\SetWatermarkScale{1}
\SetWatermarkLightness{0.99}

%------------- HEADER & FOOTER (for page 2 onward) -------------
\fancypagestyle{main}{
\fancyhf{}
\fancyhead[L]{\CourseCode\ \CourseName{} -- \DocType}
\fancyhead[R]{\href{https://qrsntu.org}{QRS@NTU}}
\fancyfoot[C]{\thepage}
\renewcommand{\headrulewidth}{0.4pt}
\renewcommand{\footrulewidth}{0pt}
}

%------------- SIMPLE FOOTER (page 1 only) -------------
\fancypagestyle{plainfirst}{
\fancyhf{}
\renewcommand{\headrulewidth}{0pt}
\renewcommand{\footrulewidth}{0pt}
\fancyfoot[C]{\scriptsize\textcolor{gray}{\textit{For academic reference only.
This document is provided for educational purposes and does not constitute an
official question paper, answer key, or any formally authorised academic
material. Its content is not guaranteed to be complete or accurate.}}}
}

%------------- HYPERREF METADATA -------------
\hypersetup{
colorlinks=true,
linkcolor=blue,
filecolor=magenta,
urlcolor=cyan,
pdftitle={\CourseCode\ \CourseName{} -- \DocType},
pdfauthor={Quantitative Research Society, NTU},
pdfsubject={\CourseName},
pdfkeywords={microeconomics, QRS@NTU, revision notes}
}

%==================================================
% DOCUMENT
%==================================================
\begin{document}
\maketitle
\thispagestyle{plainfirst}

\begin{abstract}
\noindent
Comprehensive revision notes for \CourseCode\ \CourseName, integrating theoretical concepts, key formulas, worked examples, and exam-focused applications.
\end{abstract}

\vspace{0.5cm}
\tableofcontents
\newpage

\newgeometry{left=0.7in,right=0.7in,top=1in,bottom=1in}
\fontsize{10.5pt}{11.5pt}\selectfont
\pagestyle{main}

%====================================================================
\section*{Course Overview \& Topic Map}
\addcontentsline{toc}{section}{Course Overview & Topic Map}
%====================================================================

\begin{center}
\renewcommand{\arraystretch}{1.3}
\setlength{\tabcolsep}{6pt}

\begin{tabularx}{\textwidth}{
|>{\raggedright\arraybackslash}p{3.2cm}
|>{\raggedright\arraybackslash}p{1.6cm}
|>{\raggedright\arraybackslash}X
|>{\raggedright\arraybackslash}p{1.8cm}|
}
\hline
\textbf{Topic Area} & \textbf{Study Weeks} & \textbf{Key Concepts} & \textbf{ILOs} \\
\hline
Goods Market (closed economy) &
Week 1 (13 Jan) &
Equilibrium output and demand components; consumption/saving behaviour; government spending and taxes; multiplier intuition (placeholder). &
1,2 \\
\hline
Financial Market &
Week 2 (20 Jan) &
Money vs bonds; interest rate determination; money demand and money supply (placeholder). &
1,2 \\
\hline
IS--LM Model &
Week 3 (27 Jan) &
Joint equilibrium in goods and financial markets; IS relation; LM relation; policy shifts and comparative statics (placeholder). &
1,2,4 \\
\hline
Extended IS--LM &
Week 4 (3 Feb) &
Extensions to baseline IS--LM (e.g., risk, expectations, constraints); richer policy transmission channels (placeholder). &
1,2,4 \\
\hline
Labour Market &
Week 5 (10 Feb) &
Wage setting and price setting; natural unemployment; markups; labour market equilibrium (placeholder). &
1,2 \\
\hline
Inflation and Phillips Curve &
Week 6 (17 Feb) &
Phillips curve logic; inflation dynamics; expectations; short-run trade-offs and policy implications (placeholder). &
2,4 \\
\hline
Assessment checkpoint &
Week 7 (24 Feb) &
Quiz 1 (Weeks 1--5); consolidation and error analysis (placeholder). &
1,2 \\
\hline
Recess Week &
Recess &
Catch-up, problem sets, and model integration across the short-run block (placeholder). &
1,2 \\
\hline
IS--LM--PC Framework &
Week 8 (10 Mar) &
Short-run to medium-run adjustment; interaction of output, unemployment, and inflation (placeholder). &
1,2,4 \\
\hline
Facts of Growth &
Week 9 (17 Mar) &
Stylised facts; growth accounting; capital and productivity (placeholder). &
3 \\
\hline
Economic Growth (I) &
Week 10 (24 Mar) &
Baseline growth model; capital accumulation; saving and steady state (placeholder). &
3 \\
\hline
Economic Growth (II) &
Week 11 (31 Mar) &
Technological progress; long-run growth; transitional dynamics (placeholder). &
3 \\
\hline
Economic Growth (III) \& Challenges &
Week 12 (7 Apr) &
Modern growth topics; policy debates; interpretation of evidence (placeholder). &
3,4 \\
\hline
Revision \& exam preparation &
Week 13 (14 Apr) &
Integration across business cycle and growth blocks; exam strategy and synthesis (placeholder). &
1,2,3,4 \\
\hline
\end{tabularx}


\end{center}

\noindent
“This revision guide covers all study weeks organized by economic models and policy questions rather than chronological order, ensuring comprehensive coverage while highlighting conceptual connections.”

\newpage

%====================================================================
\section*{How to Use This Document}
\addcontentsline{toc}{section}{How to Use This Document}
%====================================================================

\subsection*{Structure of the Notes}

These notes are designed to be used together with your lectures and tutorials. Each chapter is organised into:

\begin{enumerate}
\item \textbf{Core Theory} — formal definitions, results and economic intuition.
\item \textbf{Mathematical Framework} — formulas, algebra and derivations.
\item \textbf{Worked Examples} — tutorial-style questions with full solutions.
\item \textbf{Applications} — policy or real-world interpretations where relevant.
\item \textbf{Exam Focus} — high-yield concepts, typical traps and checklist items.
\end{enumerate}

A good way to revise:
\begin{itemize}
\item First pass: read Core Theory + Definition/Theorem + Key Formula boxes.
\item Second pass: attempt Worked Examples yourself \emph{before} checking solutions.
\item Final pass: review Exam Focus and Common Mistake boxes in the week before the exam.
\end{itemize}

\subsection*{Colour and Box System (Legend)}

The document uses a colour-coded system. Each box has a \textbf{name} at the top so you can quickly recognise its purpose:

\begin{center}
\begin{tabular}{|l|l|p{7cm}|}
\hline
\textbf{Colour} & \textbf{Box Name} & \textbf{Use This For} \\
\hline
\colorbox{green!5}{Green} & \textbf{Definition} & Precise concepts, variables, and model primitives to be used consistently throughout the course. \\
\hline
\colorbox{blue!3}{Blue} & \textbf{Theorem} & Core results and formal statements that structure the model logic (state assumptions and conclusions). \\
\hline
\colorbox{blue!5}{Blue} & \textbf{Key Formula} & Core equations and identities that you should memorise or be able to derive quickly. \\
\hline
\colorbox{brown!5}{Brown} & \textbf{Worked Example} & Fully worked questions. Try on your own before reading the solution. \\
\hline
\colorbox{red!5}{Red} & \textbf{Common Mistake} & Typical errors, misleading shortcuts and conceptual traps to avoid in exams. \\
\hline
\colorbox{yellow!10}{Yellow} & \textbf{Policy Application} & Real-world or policy-related interpretations to help build intuition and essay points. \\
\hline
\colorbox{purple!5}{Purple} & \textbf{Exam Focus} & High-yield checklists, ``must-know'' results and exam strategy tips. \\
\hline
\end{tabular}

\end{center}

\noindent
When you print in black and white, you can still identify boxes by their \textbf{title text}:
\begin{itemize}
\item Look for the name at the top: \textit{Definition}, \textit{Theorem}, \textit{Key Formula}, \textit{Worked Example}, \textit{Common Mistake}, \textit{Policy Application}, \textit{Exam Focus}.
\item In digital version, use the colour + name; in printed version, rely mainly on the name.
\end{itemize}


%==================================================
\part{Main Notes}
%==================================================

% ==========================================================
% HE2002 Macroeconomics II — Chapter 1: The Goods Market
% Source materials: :contentReference[oaicite:0]{index=0}
% (This is a chapter file intended to be \input into the main template.)
% ==========================================================

\section{The Goods Market}

\subsection{GDP and the Composition of Demand}

\begin{definition}{GDP (expenditure approach) and demand components}{def:gdp_components}
Gross Domestic Product (GDP) is commonly measured using the expenditure approach:
\emph{total spending on domestically produced final goods and services}.
Key components are:
consumption $C$, investment $I$, government purchases $G$, exports $X$, and imports $IM$.
\end{definition}

\begin{keyformula}[title={Key Formula: Total demand for goods}]
\[
Z \equiv C + I + G + X - IM,
\qquad
NX \equiv X-IM.
\]
\end{keyformula}

\noindent
In these notes we focus on a \textbf{closed economy} (the baseline for HE2002 short-run analysis), so $X=IM=0$ and
\[
Z \equiv C + I + G.
\]
We also \textbf{ignore inventory investment} in the model; nevertheless, inventory adjustments provide a useful
economic mechanism for how output adjusts when production differs from desired spending (see below).

\subsection{Model Setup: Behavioural Equations and Policy Variables}

\begin{definition}{Disposable income and net taxes}{def:disposable_income}
Disposable income is income available to households after taxes net of transfers:
\[
Y_D \equiv Y - T,
\]
where $Y$ is income/output and $T$ denotes taxes minus government transfers (net taxes).
\end{definition}

\begin{definition}{Consumption function (baseline Keynesian specification)}{def:consumption_function}
Consumption is assumed to depend positively on disposable income:
\[
C = C(Y_D), \qquad \frac{dC}{dY_D} > 0.
\]
A standard linear specification is
\[
C = c_0 + c_1 Y_D,
\quad 0 < c_1 < 1,
\]
where $c_1$ is the propensity to consume and $c_0$ is autonomous consumption.
\end{definition}

\begin{keyformula}[title={Key Formula: Consumption written in terms of $Y$ and $T$}]
Using $Y_D=Y-T$,
\[
C = c_0 + c_1(Y-T).
\]
\end{keyformula}

\begin{definition}{Endogenous vs.\ exogenous variables}{def:endo_exo}
\textbf{Endogenous} variables are determined within the model (e.g.\ equilibrium output $Y$).
\textbf{Exogenous} variables are taken as given (e.g.\ policy instruments or parameters).
A bar (e.g.\ $\bar I$) indicates an exogenous constant.
\end{definition}

\noindent
For the baseline goods-market model:
\begin{itemize}
\item Investment is taken as exogenous: $I=\bar I$ (this assumption is later relaxed in IS--LM analysis).
\item Fiscal policy variables $G$ and $T$ are taken as exogenous choices of the government.
\item We ignore the interest-rate effect on consumption in this chapter.
\end{itemize}

\subsection{Goods-Market Equilibrium and Equilibrium Output}

\begin{keyformula}[title={Key Formula: Planned expenditure in a closed economy}]
With $I=\bar I$,
\[
Z = C + \bar I + G
  = \bigl[c_0 + \bar I + G - c_1T\bigr] + c_1Y.
\]
\end{keyformula}

\begin{theorem}{Equilibrium output in the baseline goods-market model}{equilibrium_output}
Goods-market equilibrium requires output to equal demand:
\[
Y = Z.
\]
Under the linear consumption function and $I=\bar I$, equilibrium output is
\[
Y
= \frac{1}{1-c_1}\Bigl[c_0 + \bar I + G - c_1T\Bigr],
\qquad 0<c_1<1.
\]
\end{theorem}

\noindent
\textbf{Derivation.} Start from the equilibrium condition $Y=Z$ and substitute for $Z$:
\[
Y = c_0 + c_1(Y-T) + \bar I + G
  = c_0 + c_1Y - c_1T + \bar I + G.
\]
Collect $Y$ terms on the left:
\[
(1-c_1)Y = c_0 + \bar I + G - c_1T.
\]
Divide by $(1-c_1)$ (positive since $0<c_1<1$) to obtain the stated expression.

\begin{definition}{Autonomous spending}{def:autonomous_spending}
In the baseline model, define \emph{autonomous spending} as the component of planned expenditure not
depending on output:
\[
A \equiv c_0 + \bar I + G - c_1T.
\]
Then equilibrium output can be written compactly as $Y=\dfrac{1}{1-c_1}A$.
\end{definition}

\begin{commonmistake}
Do not confuse:
\begin{itemize}
\item the \textbf{accounting identity} defining demand ($Z \equiv C+I+G$ in a closed economy), and
\item the \textbf{equilibrium condition} ($Y=Z$), which pins down equilibrium output.
\end{itemize}
The identity holds by definition; the equilibrium condition is an economic restriction that holds when the goods market clears.
\end{commonmistake}

\subsection{Adjustment Mechanism and the Keynesian Cross Intuition}

\noindent
Even though we ignore inventory investment in the model, unintended inventories provide intuition for why output
moves toward equilibrium:
\begin{itemize}
\item If $Y<Z$, desired spending exceeds production; inventories fall unexpectedly and firms raise production.
\item If $Y>Z$, production exceeds desired spending; inventories accumulate unexpectedly and firms cut production.
\end{itemize}

\noindent
Graphically, the Keynesian cross plots $Z$ against $Y$ together with the $45^\circ$ line (where $Z=Y$).
The intersection determines equilibrium output.

\begin{center}
\begin{tikzpicture}[scale=0.95]
  \draw[->] (0,0) -- (8,0) node[below] {$Y$};
  \draw[->] (0,0) -- (0,6) node[left] {$Z$};

  % 45-degree line
  \draw (0,0) -- (7.2,5.4) node[above right] {$Z=Y$};

  % Demand line Z = A + c1 Y (illustrative)
  \draw (0,1.2) -- (7.2,4.8) node[above right] {$Z = A + c_1Y$};

  % Mark an equilibrium point (illustrative)
  \draw[dashed] (4.8,0) -- (4.8,3.6);
  \draw[dashed] (0,3.6) -- (4.8,3.6);
  \fill (4.8,3.6) circle (2pt) node[above left] {$Y^\ast$};
\end{tikzpicture}
\end{center}

\subsection{Fiscal Multipliers in the Baseline Model}

\begin{keyformula}[title={Key Formula: Multipliers (holding other exogenous variables fixed)}]
From
\(
Y=\dfrac{1}{1-c_1}\bigl[c_0+\bar I+G-c_1T\bigr],
\)
\[
\frac{\partial Y}{\partial G} = \frac{1}{1-c_1},
\qquad
\frac{\partial Y}{\partial T} = -\frac{c_1}{1-c_1}.
\]
\end{keyformula}

\noindent
\textbf{Interpretation.}
Government spending enters demand one-for-one, while taxes affect demand only through consumption.
A one-unit increase in $T$ reduces disposable income by one unit, but consumption falls only by $c_1$ units.

\begin{keyformula}[title={Key Formula: Geometric-series intuition for the multiplier}]
If autonomous spending increases by $\Delta A$, then output rises in successive rounds:
\[
\Delta Y
= \Delta A + c_1\Delta A + c_1^2\Delta A + \cdots
= \frac{1}{1-c_1}\Delta A.
\]
\end{keyformula}

\begin{examplebox}[title={Worked Example (Tutorial 1, Question 1): Equilibrium GDP, disposable income, consumption}]
\textbf{Given} (billions of euros):
\[
C = 480 + 0.5Y_D,\quad I=110,\quad T=70,\quad G=250.
\]
Solve for:
\begin{enumerate}[label=(\alph*)]
\item equilibrium output $Y$,
\item disposable income $Y_D$,
\item consumption $C$.
\end{enumerate}

\textbf{Solution.}
Disposable income is
\[
Y_D = Y - T = Y - 70.
\]
Substitute into the consumption function:
\[
C = 480 + 0.5(Y-70)
  = 480 + 0.5Y - 35
  = 445 + 0.5Y.
\]
Total demand (closed economy) is
\[
Z = C + I + G
  = (445 + 0.5Y) + 110 + 250
  = 805 + 0.5Y.
\]
Equilibrium requires $Y=Z$:
\[
Y = 805 + 0.5Y
\;\Rightarrow\;
0.5Y = 805
\;\Rightarrow\;
Y = 1610.
\]
Then
\[
Y_D = Y - 70 = 1610 - 70 = 1540,
\qquad
C = 480 + 0.5\cdot 1540 = 480 + 770 = 1250.
\]
\end{examplebox}


\newpage
\subsection{Balanced Budget Multiplier}

\begin{examplebox}[title={Worked Example (Tutorial 1, Question 3): Balanced budget multiplier}]
Start from the equilibrium expression
\[
Y = \frac{1}{1-c_1}\bigl[c_0 + \bar I + G - c_1T\bigr],
\qquad 0<c_1<1.
\]
\begin{enumerate}[label=(\alph*)]
\item By how much does $Y$ increase when $G$ increases by one unit?
\item By how much does $Y$ change when $T$ increases by one unit?
\item Why are the answers to (a) and (b) different?
\item Suppose initially $G=T$ and policy changes keep the budget balanced: $\Delta G=\Delta T=1$.
What is $\Delta Y$? Are balanced-budget changes neutral?
\item How does the value of $c_1$ affect your answer to (d)?
\end{enumerate}

\textbf{Solution.}
\begin{enumerate}[label=(\alph*)]
\item Differentiate with respect to $G$:
\[
\frac{\partial Y}{\partial G} = \frac{1}{1-c_1}.
\]
So a one-unit increase in $G$ raises $Y$ by $\dfrac{1}{1-c_1}$.

\item Differentiate with respect to $T$:
\[
\frac{\partial Y}{\partial T} = -\frac{c_1}{1-c_1}.
\]
So a one-unit increase in $T$ reduces $Y$ by $\dfrac{c_1}{1-c_1}$.

\item A change in $G$ shifts demand directly (one-for-one).
A change in $T$ affects demand only through consumption via disposable income, so only the fraction $c_1$ feeds into spending.

\item With $\Delta G=\Delta T=1$,
\[
\Delta Y
= \frac{1}{1-c_1}\Delta G - \frac{c_1}{1-c_1}\Delta T
= \frac{1}{1-c_1} - \frac{c_1}{1-c_1}
= 1.
\]
Thus a balanced-budget increase in $(G,T)$ is \textbf{not} neutral: output rises by one unit.

\item The balanced-budget multiplier is
\(
\Delta Y = 1
\)
for any $0<c_1<1$.
While $\dfrac{\partial Y}{\partial G}$ and $\dfrac{\partial Y}{\partial T}$ depend on $c_1$, their balanced-budget net effect does not.
\end{enumerate}
\end{examplebox}



\newpage
\subsection{Automatic Stabilizers: Income-Dependent Taxes}

\begin{definition}{Automatic stabilizers (in the goods-market model)}{def:automatic_stabilizers}
Taxes often increase with income. If net taxes follow
\(
T=t_0+t_1Y
\)
with $0<t_1<1$, then part of any increase in income is automatically taxed away,
dampening fluctuations in disposable income, consumption, and output.
\end{definition}

\begin{examplebox}[title={Worked Example (Tutorial 1, Question 4): Income-dependent taxes and the multiplier}]
Consider:
\[
C = c_0 + c_1Y_D,\qquad
T = t_0 + t_1Y,\qquad
Y_D = Y - T,
\]
with constant $G$ and $I=\bar I$, and $0<t_1<1$.
\begin{enumerate}[label=(\alph*)]
\item Solve for equilibrium output $Y$.
\item What is the multiplier on autonomous spending? Compare $t_1=0$ vs.\ $t_1>0$.
\item Explain why this tax system acts as an automatic stabilizer.
\end{enumerate}

\textbf{Solution.}
First compute disposable income:
\[
Y_D = Y - (t_0+t_1Y) = (1-t_1)Y - t_0.
\]
Then consumption is
\[
C = c_0 + c_1Y_D
  = c_0 + c_1\bigl[(1-t_1)Y - t_0\bigr]
  = (c_0 - c_1t_0) + c_1(1-t_1)Y.
\]

\begin{enumerate}[label=(\alph*)]
\item Goods-market equilibrium $Y=C+\bar I+G$ implies
\[
Y = (c_0 - c_1t_0) + c_1(1-t_1)Y + \bar I + G.
\]
Move the $Y$ term to the left:
\[
\bigl[1-c_1(1-t_1)\bigr]Y = (c_0 - c_1t_0) + \bar I + G.
\]
Hence
\[
Y = \frac{(c_0 - c_1t_0) + \bar I + G}{1-c_1(1-t_1)}.
\]

\item The multiplier on autonomous spending (the numerator term) is
\[
\frac{1}{1-c_1(1-t_1)}.
\]
If $t_1=0$, it reduces to $\dfrac{1}{1-c_1}$.
If $t_1>0$, then $1-c_1(1-t_1) > 1-c_1$, so the multiplier is smaller:
income-dependent taxes dampen the response of $Y$ to shocks.

\item When $Y$ rises, $T$ rises automatically, reducing $Y_D$ and hence moderating the increase in $C$.
When $Y$ falls, $T$ falls automatically, cushioning $Y_D$ and $C$.
This automatic offset stabilizes output without requiring discretionary policy changes.
\end{enumerate}
\end{examplebox}

\begin{policynote}
\textbf{Policy implication (short run).}
In the baseline goods-market framework, discretionary fiscal expansions (higher $G$ or lower $T$) increase equilibrium output through the multiplier.
With income-dependent taxes, the economy becomes less sensitive to shocks: automatic stabilizers reduce the effective multiplier, dampening both booms and recessions.
\end{policynote}

\subsection{Investment Equals Saving: An Equivalent Equilibrium Condition}

\begin{definition}{Private, public, and national saving}{def:saving_definitions}
Private saving is disposable income not consumed:
\[
S \equiv Y_D - C = Y - T - C.
\]
Public saving is the government budget surplus:
\[
T-G \quad \bigl(>0\ \text{surplus},\ <0\ \text{deficit}\bigr).
\]
National saving is the sum of private and public saving: $S + (T-G)$.
\end{definition}

\begin{keyformula}[title={Key Formula: Investment equals national saving (closed economy)}]
Starting from goods-market equilibrium $Y=C+I+G$,
\[
Y - T - C = I + G - T.
\]
Since $S \equiv Y-T-C$, this becomes
\[
S = I + G - T
\quad\Longleftrightarrow\quad
I = S + (T-G).
\]
\end{keyformula}

\begin{commonmistake}
In a closed economy with a government, the correct saving--investment relation is
\[
I = S + (T-G),
\]
not $I=S$.
Forgetting the public saving term $(T-G)$ leads to incorrect conclusions about how fiscal deficits relate to investment.
\end{commonmistake}

\noindent
\textbf{Deriving the same equilibrium output via saving.}
From $C=c_0+c_1(Y-T)$, private saving is
\[
S = Y - T - C
  = Y - T - c_0 - c_1(Y-T)
  = -c_0 + (1-c_1)(Y-T).
\]
Impose the equilibrium condition $I = S + (T-G)$ with $I=\bar I$:
\[
\bar I = \bigl[-c_0 + (1-c_1)(Y-T)\bigr] + (T-G).
\]
Expand and collect $Y$ terms:
\[
\bar I = -c_0 + (1-c_1)Y - (1-c_1)T + T - G
      = -c_0 + (1-c_1)Y + c_1T - G.
\]
Move constants to the right:
\[
(1-c_1)Y = c_0 + \bar I + G - c_1T
\quad\Rightarrow\quad
Y = \frac{1}{1-c_1}\bigl[c_0+\bar I+G-c_1T\bigr],
\]
which matches Theorem~\ref{thm:equilibrium_output}.

\begin{examtip}
High-frequency exam tasks for the goods-market chapter:
\begin{itemize}
\item Write down the model primitives and classify variables as endogenous vs.\ exogenous.
\item Derive equilibrium output algebraically from $Y=Z$; simplify carefully.
\item Compute $\frac{\partial Y}{\partial G}$ and $\frac{\partial Y}{\partial T}$; explain intuitively why the tax multiplier differs.
\item Solve for $Y$ under a balanced-budget change $\Delta G=\Delta T$ and interpret the result.
\item Use the saving framework to show $I=S+(T-G)$ and derive equilibrium output as a consistency check.
\item Explain (in words) how income-dependent taxes reduce the multiplier and why this is an automatic stabilizer.
\end{itemize}
\end{examtip}


\appendix

%==================================================
\section{Formula & Key Results Reference}
%==================================================

\noindent\textbf{Scaffold only (to be filled).} Organise key equations/results by topic:

\begin{itemize}
\item \textbf{Goods Market / IS foundations}
\begin{itemize}
\item Equilibrium condition (placeholder): \textit{planned spending} $=$ \textit{output}.
\item Multiplier and comparative statics (placeholder bullets).
\item IS relation statement (placeholder).
\end{itemize}


\item \textbf{Money \& Financial Market / LM foundations}
\begin{itemize}
    \item Money demand function and real balances (placeholder).
    \item LM relation and shifts (placeholder).
\end{itemize}

\item \textbf{IS--LM equilibrium and policy}
\begin{itemize}
    \item Joint equilibrium system (placeholder).
    \item Fiscal vs monetary policy shifts (placeholder).
    \item Extended IS--LM elements (placeholder).
\end{itemize}

\item \textbf{Labour Market}
\begin{itemize}
    \item Wage-setting and price-setting relations (placeholder).
    \item Natural rate and markups (placeholder).
\end{itemize}

\item \textbf{Inflation / Phillips Curve}
\begin{itemize}
    \item Phillips curve (expectations-augmented) (placeholder).
    \item Disinflation and sacrifice ratio intuition (placeholder).
\end{itemize}

\item \textbf{IS--LM--PC (short to medium run)}
\begin{itemize}
    \item Dynamic adjustment logic (placeholder).
    \item Medium-run equilibrium (placeholder).
\end{itemize}

\item \textbf{Growth}
\begin{itemize}
    \item Growth accounting identity (placeholder).
    \item Solow model steady state and transitional dynamics (placeholder).
    \item Technology and long-run growth mechanism (placeholder).
\end{itemize}


\end{itemize}

%==================================================
\section{Notation}
%==================================================

\noindent\textbf{Scaffold only (to be filled).}

\begin{center}
\begin{tabular}{|p{2.2cm}|p{7.8cm}|p{5.0cm}|}
\hline
\textbf{Symbol} & \textbf{Meaning} & \textbf{Context} \
\hline
$Y$ & Output / income (placeholder) & Goods market, IS--LM \
\hline
$C$ & Consumption (placeholder) & Goods market \
\hline
$I$ & Investment (placeholder) & Goods market, IS relation \
\hline
$G$ & Government spending (placeholder) & Fiscal policy \
\hline
$T$ & Taxes / net taxes (placeholder) & Fiscal policy \
\hline
$i$ or $r$ & Nominal / real interest rate (placeholder) & Financial market, LM \
\hline
$M$ & Nominal money supply (placeholder) & Financial market \
\hline
$P$ & Price level (placeholder) & Inflation, real balances \
\hline
$u$ & Unemployment rate (placeholder) & Labour market, Phillips curve \
\hline
$\pi$ & Inflation rate (placeholder) & Phillips curve, policy \
\hline
$A$ & Technology / productivity (placeholder) & Growth \
\hline
$K$ & Capital stock (placeholder) & Growth \
\hline
\end{tabular}
\end{center}

%==================================================
\section{Exam & Admin Information}
%==================================================

\noindent\textbf{Scaffold only (to be filled).} Keep factual items short.

\begin{itemize}
\item \textbf{Assessment breakdown (placeholder for policy details):}
\begin{itemize}
\item Quiz 1: 40%
\item Course Participation: 10%
\item Final Examination: 50%
\end{itemize}

```
\item \textbf{Key dates (placeholders; verify on NTULearn / course announcements):}
\begin{itemize}
    \item Quiz 1: Tuesday, 24 Feb, 9:30--11:00 (venue per course outline).
    \item Final Examination: Wednesday, 29 Apr, 5:00--7:00.
\end{itemize}

\item \textbf{Tutorials (placeholder):}
\begin{itemize}
    \item Tutorials commence in Week 2; participation recorded each tutorial (details to be inserted).
\end{itemize}

\item \textbf{Make-up quiz/test (placeholder summary):}
\begin{itemize}
    \item Eligibility and documentation requirements (placeholder).
    \item Request timeline and scheduling window (placeholder).
\end{itemize}

\item \textbf{Exam checklist (placeholder):}
\begin{itemize}
    \item Know each model: assumptions, equilibrium conditions, graphs, and policy transmission.
    \item Drill: comparative statics, diagram shifts, and short explanations in words.
    \item Maintain an error log from tutorials/quizzes and revisit weekly.
\end{itemize}


\end{itemize}

\end{document}
===== END FILE: 